\documentclass[12pt]{article}

\usepackage{exsheets}

%\usepackage[pdf]{pstricks}
%\usepackage{pstricks-add}

%\usepackage{tikz} 
\usepackage[fleqn]{amsmath}
\usepackage{amssymb}
\usepackage{amstext}  
\usepackage{amsthm} 
\usepackage{array}
\usepackage{mathtools}
\mathtoolsset{showonlyrefs} 
\usepackage{graphicx,hyperref}
\usepackage{listings}
\usepackage{indentfirst}
\usepackage[top=1.5cm, bottom=1.5cm, left=.5in, right=.5in]{geometry}
\usepackage{enumerate}
\usepackage{multicol}

\newcommand{\ssol}{\vspace{3em}}
\newcommand{\lsol}{\vspace{10em}}
\newcommand{\blnk}{{\underline {\hspace{1.5in}}}}

%\usepackage{mcode}

\setcounter{secnumdepth}{0} % Turns off automatic numbering of sections.

\SetupExSheets{headings=block}



\begin{document}
%Herbs help on direction of project
\emph{UPDATE 21-June-2016}
Per Herb's email:
\begin{quote}
Yes you can map polynomials with integer coefficients to the integers in a meaningful way by mapping 
$X \rightarrow 10$ and then showing that there is only one extension of that map that preserves addition and multiplication.

So a general project around how are the integers and the polynomials with integer coefficients similar and how are they different.  Ditto for the ring of polynomials with inter coefficients and the ring of polynomials with rational coefficients.  You could compare, for example, which of these have unique factorization, in which are all ideals principal, etc.  You could investigate how you do clock arithmetic in these systems like clock arithmetic works with respect to the integers.  You could also investigate their use in coding theory.
\end{quote}

\mbox{}\\

Justin Boyer
Possible project idea -- Striking similarities between integers and integer valued polynomials\\

There are already a lot of similarities that are shared between the integers and integer valued polynomials.  Such has:
\begin{itemize}
\item Non-zero integers admit unique factorization into primes and non-zero polynomials admit unique factorization into irreducibles.
\item The integers modulo the ideal generated by any prime form a field and the polynomials over a field modulo the ideal generated by any irreducible polynomial form a field
\item Any ideal in either of the rings is principal.
\item Euclid's algorithm can find the gcd of two integers and likewise with polynomials.
\item Both are closed under addition, subtraction and multiplication, but not under division.
\end{itemize}


What I seek to do is boil these similarities down to a common connection.  I believe this common connection is the marking of different components (as in vectors).  I intend to produce (if possible) or analyze a map that takes  integers to integer valued polynomials.  I will analyze the inverse of this map and what domains are admitted to create an isomorphism. If time permits this same analysis could be expanded to include vectors as well.  If possible I anticipate that the expansion into vectors would nicely and succinctly illustrate why polynomials, integers and vectors share so many similarities.  

I believe that this will be of benefit to teachers of mathematics because it sheds light on mathematics that is never really explained in a comprehensive manner to students.  I think this might be why a lot of "western" students struggle with numeracy.  I might be mistaken but some "eastern" countries supposedly use an additive system plus the abacus.  It also gives them a way to see the connections among seemingly unrelated objects.  Also I anticipate there could be fun math at a high school level to analyze, a few of which are: 
\begin{itemize}
\item $x^2+2x+1$ is the same as $121$ in base ten or $1b^2+2b^1+1b^0$, what do the roots of this polynomial tell us?  That 121 is 0 in base $-1$.  What does this mean?  Can we apply this notion to polynomials, if so how does this notion apply to polynomials?  What about $<1, 2, 1>$?

\item What would happen if we didn't have a positional number system?  What alternatives are there?  Would these be more or less similar to the polynomials? 

\end{itemize}



\end{document}


\begin{pspicture}
\psgrid[unit=.05\linewidth, subgriddiv=1, linewidth=.01 gridlabels=0pt, gridcolor=gray, griddots=10](0,0)(-10,-10)(10,10) 
\psaxes[unit=.05\linewidth, linewidth=.1 labelFontSize=\scriptscriptstyle]{<->}(0,0)(-10.5,-10.5)(10.5,10.5)
\end{pspicture}


\begin{tikzpicture}[scale=.5]	
\draw[help lines] (0,0) grid (20,20);
\end{tikzpicture}


